% Related Work
\section{Related Work}
\label{sec:related_work}

\subsection{Autonomous-Driving Datasets and Benchmark Bias}
\label{sec:related_datasets}

Progress in autonomous-driving perception has been strongly shaped by public benchmarks.
The KITTI Vision Benchmark Suite established one of the earliest widely adopted
real-world evaluation platforms for road-scene understanding by combining synchronized
cameras, LiDAR, and vehicle-localization sensors~\cite{geiger2012kitti}. Its
object-detection benchmark enabled reproducible comparison of cars, pedestrians, and
cyclists and remains influential in the field. However, KITTI is comparatively compact
and reflects a limited set of geographic, traffic, camera, and environmental conditions,
so strong performance on a KITTI validation split does not by itself demonstrate
robustness outside the source distribution.

Later datasets broadened the scale and diversity of driving perception. The Waymo Open
Dataset provides large-scale camera and LiDAR recordings collected across multiple cities
and driving environments~\cite{sun2020waymo}. BDD100K contributes extensive geographic,
weather, illumination, and scene variation across a large collection of driving videos and
multiple perception tasks~\cite{yu2020bdd100k}. nuScenes further expands benchmark
diversity through a full sensor suite, 360-degree coverage, radar data, and dense
multimodal annotations~\cite{caesar2020nuscenes}. Together, these datasets show that
autonomous-driving benchmarks differ not only in size but also in viewpoint, image
characteristics, traffic density, object-size distribution, class frequency, annotation
ontology, and operating environment.

These differences create a meaningful benchmark-bias problem. A detector can exploit
regularities specific to its training dataset---camera height, road layout, background
appearance, and class prevalence---while appearing highly accurate under an in-domain
split. Cross-dataset testing is therefore a necessary complement to conventional
validation. Nevertheless, dataset papers mainly define benchmarks and evaluation tasks;
they do not provide a controlled comparison in which detectors trained under the same
protocol on KITTI are transferred directly to a harmonized Waymo subset. This limitation
motivates the external-validation setting adopted here.

\subsection{CNN-Based Object Detection}
\label{sec:related_cnn}

Convolutional object detectors are commonly organized into two-stage and one-stage
paradigms. Faster R-CNN integrated a Region Proposal Network with a region-based detector,
allowing proposal generation and classification to share convolutional
features~\cite{ren2015fasterrcnn}. This design established a strong accuracy-oriented
baseline and remains relevant when localization quality is prioritized, although its
proposal and per-region processing stages can increase computational cost relative to
compact one-stage architectures.

RetinaNet addressed a central limitation of dense one-stage detection: the extreme
imbalance between background locations and foreground objects. Focal loss reduces the
contribution of easy negatives and concentrates learning on difficult
examples~\cite{lin2017focalloss}, allowing a one-stage detector to achieve competitive
accuracy with a simpler dense prediction pipeline. This formulation is particularly
relevant to road scenes, where the image contains large background regions and
comparatively few small foreground instances. The original work, however, focused on
standard in-domain benchmarks and did not evaluate how focal-loss-based detection behaves
when transferred between autonomous-driving datasets.

The YOLO family represents the real-time one-stage paradigm; YOLOv7, for example,
introduced architectural and training refinements designed to improve accuracy without
proportionally increasing inference cost~\cite{wang2023yolov7}. Such models are
attractive for autonomous vehicles because of their high throughput and compact
deployment footprint. Published comparisons among YOLO, Faster R-CNN, and RetinaNet are
difficult to interpret directly because studies frequently differ in datasets, input
resolutions, backbones, pretraining sources, hardware platforms, precision modes, and
latency definitions. A controlled same-data and same-hardware evaluation is therefore
required to determine whether speed-oriented CNN detectors preserve their advantages
under cross-dataset shift.

\subsection{Transformer-Based and End-to-End Detection}
\label{sec:related_transformer}

Transformer-based detectors introduced a substantially different formulation of object
detection. DETR treated detection as direct set prediction and used bipartite matching to
associate predicted objects with ground-truth instances~\cite{carion2020detr}. By
combining a convolutional backbone with a transformer encoder--decoder, DETR removed
several hand-designed components, including anchor generation and non-maximum suppression.
This end-to-end formulation simplified the pipeline, but the original model required long
training schedules and was less effective for small objects, which are common and
safety-relevant in driving scenes.

Deformable DETR improved the efficiency of transformer attention by sampling a sparse set
of relevant spatial locations across multi-scale feature maps~\cite{zhu2021deformabledetr},
accelerating convergence and improving small-object performance. RT-DETR extended this
direction by explicitly targeting real-time end-to-end detection through an efficient
hybrid encoder and flexible speed control~\cite{zhao2024rtdetr}, providing an appropriate
modern transformer baseline for comparison with real-time CNN detectors.

Despite their architectural differences, CNN and transformer detectors are commonly
compared only on the benchmark on which each model was developed. It remains uncertain
whether global or selective attention, end-to-end set prediction, region proposals, and
dense convolutional prediction produce different robustness patterns under natural
autonomous-driving domain shift. This question is addressed by evaluating RT-DETR and
representative CNN families under one harmonized KITTI-to-Waymo protocol.

\subsection{Domain Shift and Cross-Dataset Generalization}
\label{sec:related_domain_shift}

Domain shift arises when the distribution encountered during deployment differs from the
distribution used for model development. In object detection this shift may involve
appearance, weather, illumination, camera geometry, object scale, traffic density, class
balance, or annotation conventions. Domain-adaptation research attempts to reduce this
mismatch by using information from the target domain: Domain Adaptive Faster R-CNN
introduced image-level and instance-level adversarial alignment for an unlabeled target
domain~\cite{chen2018domainadaptive}, and the Cross-Domain Adaptive Teacher later used
teacher--student learning with target-domain pseudo-labels~\cite{li2022crossdomainadaptive}.
These methods are valuable when target images are available during training, but they
answer a different question from target-free external validation.

Direct cross-dataset studies provide evidence that benchmark-specific performance may not
generalize. Hasan et al.\ showed that pedestrian detectors with strong within-dataset
results can degrade substantially when evaluated on other pedestrian
datasets~\cite{hasan2021generalizablepedestrian}, exposing the risk of dataset-specific
design choices, although the principal focus is pedestrian detection rather than a unified
comparison across detector paradigms. The SHIFT dataset enables controlled investigation
of continuous changes in weather, time of day, and traffic conditions~\cite{sun2022shift},
while weather-robust detection studies analyze real and synthetic adverse-weather
transformations~\cite{gupta2024weather}. Such work clarifies individual sources of
degradation but does not fully represent the combined geographic, camera, ontology, and
scene-distribution differences between independent real-world datasets.

The present study therefore distinguishes three evaluation settings. \emph{Domain
adaptation} modifies training using labeled or unlabeled target-domain information.
\emph{Domain generalization} typically modifies training to improve performance on
unspecified unseen domains. \emph{Target-free external validation}, which is used here,
leaves the trained detector unchanged and measures its direct transfer to an independently
constructed target dataset. This setting reveals the natural external reliability of
standard detectors before adaptation is introduced and supports a fair analysis of
architecture-dependent degradation.

\subsection{Safety-Oriented and Deployment-Aware Evaluation}
\label{sec:related_safety}

Mean Average Precision is indispensable but insufficient for safety-oriented
autonomous-driving evaluation. Similar mAP values can arise from different error profiles,
including misclassification, poor localization, duplicate detections, background false
positives, and missed objects. TIDE formalized this distinction through a model-agnostic
decomposition of object-detection errors~\cite{bolya2020tide}. Such analysis is
particularly important for pedestrians and cyclists, because a moderate aggregate score
can conceal a high number of safety-critical false negatives; object size, distance,
occlusion, crowding, and class frequency should therefore be examined alongside overall
accuracy.

Deployment constraints introduce an additional evaluation dimension. Streaming-perception
work demonstrated that processing delay changes the effective quality of online
predictions because a delayed result may describe an earlier scene
state~\cite{yang2022streaming}, and the ASAP benchmark similarly incorporated latency into
vision-centric autonomous-driving evaluation, showing that strong offline models may behave
differently under streaming conditions~\cite{wang2023asap}. These findings motivate
measuring detector performance on common hardware using batch-one latency, throughput,
and, where possible, latency percentiles rather than reporting only theoretical speed. A
deployment-oriented comparison should also report parameter count, computational
complexity, model size, and memory requirements.

\subsection{Research Gap and Positioning}
\label{sec:related_gap}

The reviewed literature provides strong foundations in autonomous-driving datasets, CNN
and transformer detection, domain adaptation, error analysis, and latency-aware
evaluation, but these streams remain fragmented. Dataset studies emphasize scale and
diversity; detector papers prioritize in-domain architectural performance; adaptation
methods use target-domain information; cross-dataset studies often focus on one class or
one shift type; and deployment studies usually examine latency separately from external
robustness. There is limited evidence from a controlled target-free experiment in which
proposal-based, dense one-stage, real-time CNN, and real-time transformer detectors are
trained on the same source dataset and evaluated on a distinct real-world dataset using
harmonized classes and common implementation rules. The proposed study addresses this gap
by training YOLO, Faster R-CNN, RetinaNet, and RT-DETR on KITTI and externally validating
them on a representative Waymo front-camera subset without retraining or adaptation,
jointly considering in-domain accuracy, absolute and relative degradation, class-wise
generalization, false negatives, difficult-object conditions, and same-hardware
efficiency.
